\section*{Next Steps and Planned Progress}

Multi-process isolation and range-based dirty tracking (the \texttt{dirty\_range} procfs interface) are fully implemented.
The test suite covers table lifecycle, ordering guarantees, address range filtering, and
fault-to-visibility latency. The remaining work is as follows.

\subsection*{Extensions Still Pending}

\begin{enumerate}
      \item \textbf{CPU-side Dirty Tracking:}
            The current implementation intercepts GPU write faults only. Extending tracking to
            CPU-side writes requires integration with the \texttt{mmu\_notifier} invalidation
            path so that CPU dirty pages are recorded in the same per-process table and exposed
            through the existing procfs interface.

      \item \textbf{Per-page Access Frequency Counters:}
            The tracker currently operates under a first-write-wins policy: each page is recorded
            once with its earliest observed write timestamp. Richer statistics, such as per-page
            write counts and dirty page age, are needed to support informed migration and eviction
            decisions by user-space tools. This requires extending \texttt{struct dirty\_page\_info}
            with a counter field and updating the fault hook to increment it atomically.

      \item \textbf{Moving from Invalidation to Write Permission Removal:}
            The current implementation simply invalidates the Page Table entries as a Naive implementation,
            however it would be more efficient to just remove the write permission of only the tracked pages,
            reducing overhead for non-tracked pages and for tracked read-only pages
      \item \textbf{Batched Write-Permission Revocation:}
            The current implementation invalidates all the Page Table entries start of each epoch, 
            by iterating over all the \texttt{va\_blocks} one by one 
            A dedicated kernel thread can amortise this cost by
            processing pages in batches, reducing the latency of the \texttt{start\_track}
          operation under large tracked sets.
\end{enumerate}

\subsection*{Testing and Performance Analysis}

\begin{enumerate}

    \item \textbf{Memory Footprint Profiling:}
          Under workloads with large dirty sets, the two-level xarray structure may incur
          non-trivial memory overhead. The footprint of the tracking data structures will be
          measured and compared against alternatives such as bitmaps.

    \item \textbf{Concurrency Stress Testing:}
          The existing concurrent-stream test (\texttt{tc06\_concurrent\_streams}) exercises
          the xarray under parallel GPU faults. Additional stress tests under multi-process
          workloads and simultaneous range-filter updates are needed to validate the spinlock
          boundary.
\end{enumerate}

\subsection*{Optimisation}

Based on profiling results, targeted optimisations will be applied. Likely candidates include
batching xarray inserts to reduce per-fault locking overhead, and replacing the xarray with a
bitmap representation for workloads where the dirty set density is high enough to justify the
fixed memory cost.

\subsection*{Further Deliverables}

\begin{enumerate}
    \item \textbf{Documentation.}
          Comprehensive usage documentation covering the procfs interface, epoch semantics,
          range filtering, and known limitations.

    \item \textbf{Formal Benchmarking Report.}
          Quantified results for fault latency overhead, memory footprint, and migration
          throughput impact across representative workloads, with and without the tracker active.

    \item \textbf{Command-line Interface.}
          A user-facing CLI tool that wraps the procfs interface, providing commands to start
          and stop tracking, query dirty pages for a given PID, set address range filters, and
          display results in a human-readable format.
\end{enumerate}
